\documentclass[11pt,a4paper]{article}
\usepackage[T1]{fontenc}
\usepackage[utf8]{inputenc}
\usepackage{lmodern}
\usepackage[margin=25mm]{geometry}
\usepackage{amsmath,amssymb,amsthm,mathtools}
\usepackage{booktabs,array,longtable}
\usepackage{microtype}
\usepackage{xcolor}
\usepackage{enumitem}
\usepackage{fancyhdr}
\usepackage{hyperref}
\usepackage{xurl}
\hypersetup{colorlinks=true,linkcolor=blue!45!black,citecolor=blue!45!black,
 urlcolor=blue!45!black,pdftitle={Integral-root carriers for Con(ZFC)},
 pdfauthor={AI-assisted research prepared for Vladimir Reshetnikov}}
\setlength{\emergencystretch}{2em}
\setlist{itemsep=3pt,topsep=5pt}
\pagestyle{fancy}
\fancyhf{}
\fancyhead[L]{\small Integral-root carriers for $\mathrm{Con}(\mathrm{ZFC})$}
\fancyhead[R]{\small 23 October 2026}
\fancyfoot[C]{\thepage}
\setlength{\headheight}{14pt}
\newtheorem{theorem}{Theorem}[section]
\newtheorem{lemma}[theorem]{Lemma}
\newtheorem{proposition}[theorem]{Proposition}
\newtheorem{corollary}[theorem]{Corollary}
\theoremstyle{remark}
\newtheorem{remark}[theorem]{Remark}
\newcommand{\Zp}{\mathbb Z_{>0}}
\newcommand{\Con}{\operatorname{Con}}
\newcommand{\ZFC}{\mathrm{ZFC}}
\newcommand{\Ptwo}{P_{15,24}}
\newcommand{\Pthree}{P_{15,26}}
\newcommand{\Pone}{P_{15,30}}
\newcommand{\ord}{\mathrm{old}}
\title{Integral-root carriers for $\Con(\ZFC)$\\[4pt]
\large A degree-24 equation in fifteen positive unknowns,\\
and a degree-629 companion in twelve}
\author{AI-assisted research prepared for Vladimir Reshetnikov\\
\small Proposed contribution to the ProveIt research collection}
\date{23 October 2026}

\begin{document}
\maketitle
\begin{abstract}
We give a fixed integer polynomial in fifteen strictly positive unknowns,
of exact total degree $24$, whose lack of positive zeros is provably equivalent
in ZFC to syntactic $\Con(\ZFC)$. Its division-free evaluator uses $90$ live
binary arithmetic operations. The comparable degree frontier in the supplied
ProveIt material is degree $44$ at $75$ operations. Two further tradeoffs have
degree/operation pairs $(26,88)$ and $(30,79)$. A twelve-positive-variable
companion has degree $629$ at $114$ operations, improving the supplied
degree-$897$, $108$-operation point.

The new reduction replaces expensive coordinate scales by divisibility
forced through existing monic square equations. A rational offset is placed
inside a square root and denominators are cleared; integrality then recovers
the missing scale without an extra unknown. Positive Pell gaps supply the
reverse construction. Reducing the auxiliary expression modulo an explicit
index equation lowers degree $26$ to $24$.

The fixed coefficient tuple, calibrated compiler, and standard Pell lemmas
are inherited with precise provenance. The new reductions are proved here
by ordinary mathematics formalizable in ZFC; they are not a newly checked
end-to-end proof-assistant development. Literal circuits, exact algebra
certificates, independent finite checks, coefficient payloads, and source
proofs accompany the paper. No global minimum, consistency proof, or new
operation-count record is asserted.
\end{abstract}

\clearpage
\tableofcontents
\newpage

\section{Result, comparison, and proof boundary}
\label{sec:scope}

The goal is a particularly simple \emph{fixed} Diophantine representation of
ZFC inconsistency. The polynomial is not a family with an unspecified
machine parameter: the four numeral ports used below are bound to explicit
finite binary integers in the archive. Every other abbreviation is a
polynomial expression in the listed unknowns.

\begin{theorem}[Principal result]
\label{thm:main}
With the coefficient tuple specified in Section~\ref{sec:coefficients}, the
polynomial $\Ptwo$ of Section~\ref{sec:formula} has fifteen positive unknowns,
exact total degree $24$, and a $90$-operation evaluator. At the ordinary
mathematical proof level of the pinned native compiler,
\[
 \ZFC\vdash
 \left[\neg\exists\mathbf x\in\Zp^{15}\;\Ptwo(\mathbf x)=0\right]
 \longleftrightarrow\Con(\ZFC).
\]
The polynomials $\Pthree$ and $\Pone$ have the same equivalence, with
degree/operation pairs $(26,88)$ and $(30,79)$. The companion $P_{12}$ has
twelve positive unknowns, degree $629$, and $114$ operations.
\end{theorem}

Here $\Con(\ZFC)$ is the syntactic assertion that the fixed effective ZFC
proof calculus has no finite refutation. The inherited K215 calibration
uses that assertion, and its ordinary internalization supplies the displayed
provability claim. This notation does not claim that a new object-language
ZFC proof file or a proof-assistant certificate has been produced.

There are several distinct notions of simplicity. Table~\ref{tab:frontier}
keeps the domain, number of supplied unknowns, total degree, and literal
operation count separate. In particular, a reduction of degree at a fixed
number of positive unknowns should not be described as a reduction of the
number of unrestricted integer unknowns.

\begin{table}[htbp]
\centering
\small
\begin{tabular}{@{}lrrrr@{}}
\toprule
Construction & Positive unknowns & Degree & Operations & Status\\
\midrule
Source 47, minimum-norm grouping & 15 & 44 & 75 & Prior\\
Source 49, two root gaps & 15 & 44 & 76 & Prior\\
Source 51, sparse product & 15 & 157 & 66 & Prior\\
$\Ptwo$ & 15 & \textbf{24} & 90 & New\\
$\Pthree$ & 15 & \textbf{26} & 88 & New\\
$\Pone$ & 15 & \textbf{30} & 79 & New\\
Source 52, lean norm & 12 & 897 & 108 & Prior\\
$P_{12}$ & 12 & \textbf{629} & 114 & New\\
\bottomrule
\end{tabular}
\caption{Bounded comparisons with the supplied source collection. All entries
are single equations; systems are not counted as single polynomials.}
\label{tab:frontier}
\end{table}

Degree $24$ is $45.45\%$ below $44$, with fifteen more arithmetic gates;
degree $30$ is $31.82\%$ below $44$, with four more gates. The twelve-variable
degree reduction is $29.88\%$, with six more gates. These points make a
substantial improvement on one coordinate of the measured frontier. They
do not dominate the $66$-operation construction or the constructions with
fewer variables.

The read-only comparison includes the incoming wave-7 archives and their
subsequent source directories 46--53. The working snapshot is ProveIt
commit \path{6d375f2d2deac0519e659bbce5b48087023c1430}; retrieval started
at \path{4b4349064a750ba6347293b55874a7f92abdd0db}. The archive records
individual source hashes. Source 47 supplies the $44/75$ grouping, and
source 49 supplies the directly related $44/76$ root-gap point,
source 51 supplies the squared-modulus quotient used below, and source 52
supplies the twelve-variable predecessor \cite{source47,source49,source51,source52}.
The broader small-variable context includes Sun's ten-integer-unknown
preprint and the work on eleven integer unknowns and Isabelle
formalization \cite{sun,bayer2026,bayer2025}. Neither result is needed for
our fifteen-positive-variable proof.

The semantic dependencies are explicit. First, a finite refutation is
equivalent to a blank K215 halting computation. Second, a native compiler
encodes that computation with the authentic four-port tuple. Third, the
ordinary Pell kernel encodes the resulting arithmetic constraints. The
new work is an existence-equivalence reduction from that established
kernel to the equations here. According to the preserved K215 status
document, the exact machine endpoint already has compiled Lean and Rocq
correctness proofs \cite{k215}. We did not rerun those builds. The remaining
formalization boundary for this contribution is the native Diophantine
compiler together with the new full equation and its internalization.

\section{The mechanism: force a scale through an integral root}
\label{sec:mechanism}

\begin{lemma}[A rational square root of an integer is integral]
\label{lem:integral}
If $m>0$ and $m^2\mid n^2$ for integers $m,n$, then $m\mid n$.
Equivalently, a rational number whose square is an integer is an integer.
\end{lemma}
\begin{proof}
For every prime $p$, $2v_p(n)\ge2v_p(m)$ implies
$v_p(n)\ge v_p(m)$. The case $n=0$ is immediate. Alternatively reduce the
rational fraction to lowest terms and compare denominator prime factors.
\end{proof}

Suppose a construction uses a scaled coordinate $Y=sq^3$ and a square
equation $\tau^2=A$. If the canonical witnesses have
$\tau>L+s$, one can replace the supplied coordinate $s$ by $Y$ and the
supplied root $\tau$ by a positive gap $r$. The polynomial
\[
 [q^3(L+r)+Y]^2=q^6A
\]
forces $q^3\mid Y$ by Lemma~\ref{lem:integral}, and then restores
$\tau=L+r+Y/q^3$. The converse uses $r=\tau-L-s>0$.

The important issue is not the elementary divisibility lemma by itself.
One must find an existing root that can carry the fractional offset, prove
that no extra roots are introduced, and prove a sufficiently large positive
margin for the converse. Otherwise a change of variables can silently
exclude all intended witnesses. Here the first Pell norm supplies one such
margin, and the main Pell norm supplies another.

The second application is sequential. After recovering $q^3\mid Y$, we
know $H\equiv3\pmod q$. A canonical congruence restriction gives
$\gcd(q,3)=1$. Therefore
\[
 (qD_0+XH)^2=q^2(1+\Delta c^2)
 \quad\Longrightarrow\quad q\mid XH
 \quad\Longrightarrow\quad q\mid X.
\]
Thus both formerly scaled coordinates can be supplied directly. The
divisibility conditions have been transferred into equations already
needed for the construction, without introducing quotient unknowns.

The final degree reduction uses a separate idea: after imposing
$R=k-hXY-1$, use the lower-degree right side inside the auxiliary norm.
This is an exact reduction modulo a retained equation, not a heuristic
approximation to $R$.

\clearpage
\section{The explicit degree-24 equation}
\label{sec:formula}

The fifteen unknowns, all in $\Zp$, are
\[
 j,F,\alpha,z,f,h,i,T,X,Y,r,\eta,\zeta,u,\gamma.
\]
The four fixed integers are $\beta,K_0,M_C,M_F$. In particular
$\beta=2^{590069}-1$, and $M_F$ denotes the already shifted source port.
Define
\begin{align*}
 J&=3j+1,&q&=\beta J+1,&C&=q-F-\alpha,\\
 E&=XY,&k&=\eta+\zeta,&c&=kY+\eta,\\
 a&=E+Y,&H&=4a+3,&\Delta&=a^2+H=(a+2)^2-1,\\
 D_0&=X+ac+\gamma H,&\ell&=EkY,&I&=ic^2,\\
 G&=(q-1)(q-F)+\alpha,&&&
\end{align*}
\begin{align*}
 R&=G(q^2-1)+(M_C+qM_F)J,\\
 \widehat R&=k-hE-1,\\
 U&=\Delta(f^2-1),\\
 V&=f^2(cT-\widehat R)-c,\\
 W&=q\bigl[K_0C+(q-F)-z(q-1)\bigr]+XC.
\end{align*}
Set the six residuals
\begin{align}
 A_0&=[q^3(\ell+r)+Y]^2-q^6[\ell(\ell+k)+1],\label{eq:first}\\
 A_1&=(qD_0+XH)^2-q^2(\Delta c^2+1),\label{eq:main}\\
 A_k&=\widehat R-R,\label{eq:index}\\
 A_t&=W-q,\label{eq:transport}\\
 A_a&=(V+u)^2-Uu(2V+u)-1,\label{eq:aux}\\
 A_s&=I^2-U.\label{eq:strong}
\end{align}
The requested fixed polynomial is
\begin{equation}
 \boxed{\Ptwo=A_0^2+A_1^2+A_k^2+A_t^2+A_a^2+A_s^2=0.}
 \label{eq:p24}
\end{equation}
This is an ordinary polynomial equation over the integers. There are no
variable exponents, divisions, radicals, bit operations, or implicit
existentially quantified registers in its definition. The named expressions
only specify a compact division-free arithmetic circuit for that polynomial.
The largest exponent in the notation is a fixed numeral.

The extra roots that could arise from allowing negative supplied unknowns
are not covered by this theorem. Section~\ref{sec:domains} gives explicit
translations when the intended domain is nonnegative or unrestricted
integers.

\clearpage
\section{The inherited kernel used by the proof}
\label{sec:interface}

We state the precise interface rather than assume an unspecified
``universal Diophantine polynomial.'' Let the old positive coordinates be
\[
 J,F,\alpha,z,f,h,i,T_{\ord},s,w,\tau,\eta,\zeta,y,\gamma_{\ord}.
\]
Use the same abbreviations as above, except
\[
 X=wq,\quad Y=sq^3,\quad
 D=X+ac+\gamma_{\ord}H,\quad
 V_{\ord}=c(T_{\ord}f-1)-Rf^2.
\]
The ordinary source factors are
\begin{align*}
 N_0&=\tau^2-\ell(\ell+k),&
 N_1&=D^2-\Delta c^2,\\
 N_a&=U(V_{\ord}^2-y^2)+y^2,&
 N_k&=k-hE-R,\\
 N_t&=(K_0+w)C+(q-F)-z(q-1),&
 N_s&=1+i^2c^4-U.
\end{align*}
The inherited ordinary proof establishes soundness of
\begin{equation}
 N_0N_1N_aN_kN_tN_s=1
 \label{eq:oldproduct}
\end{equation}
on these positive coordinates: a solution yields a finite ZFC refutation.
The proof handles the possible unit signs before recovering $N_k=N_t=1$;
it does not simply assume each factor equals $+1$. Its complete ordinary
argument is preserved as \path{upstream/ORDINARY_KERNEL.md}.

Conversely, a finite refutation gives canonical normalized witnesses with
all these factors equal to $+1$. The ordinary strong coordinate is obtained
from a normalized coordinate $i_0$ by
\begin{equation}
 i=\Delta i_0,\qquad
 f^2=1+\Delta i_0^2c^4,\qquad
 I=ic^2=\Delta i_0c^2,\qquad U=I^2.
 \label{eq:normalized}
\end{equation}
The authentic compact adapter also supplies
\begin{gather}
 q=B^N,\qquad J=1+B+\cdots+B^{N-1},\qquad
 B=2^{590069},\quad N=581197^n>1,\label{eq:canonical}\\
 R\equiv3\pmod4,\qquad
 c=\psi_{a+2}(R),\qquad
 \tau=\chi_{P_0}(n_0),\quad k=2\psi_{P_0}(n_0),\label{eq:canonicalpell}\\
 P_0=2XY^2+1,\qquad n_0=(R+1)/2\ge3,\qquad k>XY,\quad q\ge16.
 \nonumber
\end{gather}
The exponent $n$ in \eqref{eq:canonical} and the Pell index $n_0$ are
proof-only integers, not supplied unknowns. Padding lets the canonical
exponent be taken arbitrarily large as required by the compiler.

Here the Pell polynomials satisfy
\[
 \chi_A(0)=1,\quad\psi_A(0)=0,\qquad
 \chi_A(n)+\psi_A(n)\sqrt{A^2-1}
       =(A+\sqrt{A^2-1})^n.
\]
They have integer coefficients in $A$. Their elementary recurrences,
classification, strong divisibility, and index-recovery lemmas are
preserved in \path{upstream/PELL_LEMMAS.md} and the ordinary kernel proof.

The coefficient interface matters. Source 53's relaxed-rank exposition
and source 51's quotient proof were originally distributed with different
historical tuples. We use their coefficient-independent algebraic lemmas,
and use source 44's \emph{whole} compact tuple and its complete native
adapter. In particular, odd $N$ and $R\equiv3\pmod4$ come from the adapter
actually bound here, not from a borrowed numeric instance.

\section{Soundness of the degree-24 construction}
\label{sec:soundness}

Suppose a positive tuple satisfies \eqref{eq:p24}. Since squares of integers
are nonnegative, all six residuals vanish.

\subsection{Recover the two coordinate scales}

Equation \eqref{eq:first} and Lemma~\ref{lem:integral} give
\[
 q^3\mid q^3(\ell+r)+Y,\qquad q^3\mid Y.
\]
Thus
\begin{equation}
 s=Y/q^3\in\Zp,\qquad \tau=\ell+r+s\in\Zp,
 \qquad \tau^2-\ell(\ell+k)=1.
 \label{eq:restorefirst}
\end{equation}
The odd exponent $590069$ gives $B\equiv2\pmod3$ and
$\beta\equiv1\pmod3$. The supplied parameterization $J=3j+1$ consequently
gives $q\equiv2\pmod3$, without assuming any decoded form for $q$.
Moreover $q\mid Y$ gives $H=4Y(X+1)+3\equiv3\pmod q$. Hence
\[
 \gcd(q,H)=\gcd(q,3)=1.
\]
Equation \eqref{eq:main} and the same lemma yield $q\mid XH$, and therefore
\begin{equation}
 w=X/q\in\Zp,\qquad
 \gamma_{\ord}=\gamma+w\in\Zp,
 \qquad D=D_0+wH.
 \label{eq:restoremain}
\end{equation}
After division by $q^2$, \eqref{eq:main} is exactly
$D^2-\Delta c^2=1$. No prime, power-of-two, or squarefree hypothesis about
an arbitrary supplied $q$ has been used.

\subsection{Recover the index, transport, and auxiliary witnesses}

The index equation gives $R=\widehat R$, hence $N_k=1$.
Using $X=qw$, the transport numerator factors exactly as
\[
 W=q\bigl[(K_0+w)C+(q-F)-z(q-1)\bigr]=qN_t.
\]
Since $W=q$ and $q>0$, $N_t=1$.

Put $T_{\ord}=fT>0$. Then
\[
 V_{\ord}=c(T_{\ord}f-1)-Rf^2
            =f^2(cT-R)-c=V.
\]
Let $y=|V+u|$. The exact identity
\begin{equation}
 U\bigl[V^2-(V+u)^2\bigr]+(V+u)^2
       =(V+u)^2-Uu(2V+u)
 \label{eq:auxidentity}
\end{equation}
shows $N_a=1$. We must also check $y>0$. From $A_s=0$,
$U=I^2$ with $I=ic^2>1$; indeed $k\ge2$ and $c=kY+\eta\ge3$.
If $y=0$, the auxiliary norm would give $I^2V^2=1$, impossible for
integer $V$ and $I>1$. Thus the restored ordinate is positive.
The strong equation also gives $f>1$ and $N_s=1$.

We have reconstructed a positive solution of \eqref{eq:oldproduct}, in
fact with all six source factors $+1$. The inherited soundness theorem
therefore supplies a finite ZFC refutation. All divisions above occur
only after their integrality has been proved; none occurs in the polynomial.

\section{Completeness: the positive gaps really exist}
\label{sec:completeness}

Start with a finite ZFC refutation and choose the canonical normalized
witnesses from Section~\ref{sec:interface}. Retain the outer coordinates
and set $i=\Delta i_0$ as in \eqref{eq:normalized}. We construct positive
$j,r,\gamma,T,u$ explicitly.

\subsection{The cofinal congruence restriction}

Both $590069$ and $N=581197^n$ are odd. Reducing the geometric sum
in \eqref{eq:canonical} modulo $3$ gives
\[
 J\equiv 1-1+1-\cdots+1\equiv1\pmod3.
\]
As $N>1$, $J>1$. Therefore $j=(J-1)/3$ is a positive integer.
This proves that imposing $J=3j+1$ retains the canonical witnesses;
it does not claim this congruence for every old, noncanonical solution.

\subsection{The first root has room for its offset}

The old first norm is $\tau^2=\ell(\ell+k)+1$. Since
$\ell=XY^2k>k$ and $k\ge2$, it gives $\tau<\ell+k$. Consequently
\begin{equation}
 \tau-\ell=\frac{\ell k+1}{\tau+\ell}>\frac{k}{3}.
 \label{eq:firstgap}
\end{equation}
Canonical $k>XY$, $X\ge q\ge16$ now imply
\[
 \frac{k}{3}>\frac{XY}{3}>\frac{Y}{q^3}=s.
\]
Thus $r=\tau-\ell-s$ is a positive integer, and \eqref{eq:first} holds
with the new directly supplied $Y=sq^3$.

\subsection{The main root has room for its offset}

The old main root satisfies $D=\sqrt{\Delta c^2+1}>(a+1)c$.
Therefore
\begin{equation}
 \gamma_{\ord}=\frac{D-ac-X}{H}>\frac{c-X}{H}.
 \label{eq:gammabound}
\end{equation}
With $P_0=2XY^2+1$ and $n_0\ge3$, the canonical first Pell coefficient
satisfies
\[
 k=2\psi_{P_0}(n_0)\ge2\psi_{P_0}(3)=8P_0^2-2>32X^2Y^4.
\]
Hence $c>kY>32X^2Y^5$. As $X\ge16$ and $Y\ge1$, we have
$H=4XY+4Y+3\le9XY$. These bounds give $c-X>XH$, so
\[
 \gamma_{\ord}>X\ge X/q=w.
\]
The new coordinate $\gamma=\gamma_{\ord}-w$ is therefore positive.
Substituting it gives $D_0=D-wH$, and hence \eqref{eq:main}.
This strict margin is essential: merely observing that $\gamma_{\ord}$
is positive would not justify subtracting $w$.

\subsection{A squared-modulus quotient and the auxiliary gap}

The squared-modulus quotient is an established ingredient from source 51,
reproved here to make the completeness step explicit \cite{source51}.
For $r_0\ge0$ there is an integer polynomial $Q_{r_0}$ such that
\begin{gather}
 \chi_t(2r_0+1)=tQ_{r_0}(t^2),\label{eq:oddpoly}\\
 Q_{r_0}(0)=(-1)^{r_0}(2r_0+1),\qquad
 Q_{r_0}(1-A^2)=(-1)^{r_0}\psi_A(2r_0+1).
 \label{eq:oddcongruence}
\end{gather}
Indeed $Q_0(x)=1$, $Q_1(x)=4x-3$, and
\[
 Q_{r_0+1}(x)=(4x-2)Q_{r_0}(x)-Q_{r_0-1}(x).
\]
The first identity follows from the odd subsequence of the Pell recurrence.
The evaluation at $0$ follows by induction. Substitution $x=1-A^2$
negates the coefficient $4A^2-2$ of the odd $\psi_A$ subsequence; its
initial values are $1,4A^2-1$, proving the second identity.

Write $A=a+2$ and $R=2r_0+1$. Canonical $R\equiv3\pmod4$ makes $r_0$
odd. Define
\begin{equation}
 V_* = \chi_I(R)/I=Q_{r_0}(I^2),\qquad y_* = \psi_I(R).
 \label{eq:auxfresh}
\end{equation}
These are positive integers, and the Pell identity gives
\begin{equation}
 I^2V_*^2-(I^2-1)y_*^2=1.
 \label{eq:auxnormfresh}
\end{equation}
Since $I^2=\Delta(f^2-1)\equiv1-A^2\pmod{f^2}$ and $c\mid I$,
\eqref{eq:oddcongruence} gives
\[
 V_*\equiv-c\pmod{f^2},\qquad V_*\equiv-R\pmod c.
\]
Moreover \eqref{eq:normalized} gives $f^2\equiv1\pmod c$, so
$\gcd(c,f^2)=1$. It follows that
\begin{equation}
 T=\frac{V_*+c+Rf^2}{cf^2}\in\Zp.
 \label{eq:newT}
\end{equation}
For this $T$, the new formula for $V$ equals $V_*$ once the index residual
vanishes. This is a fresh auxiliary witness construction; it does not
claim that every arbitrary old quotient was divisible by $f$.

Finally $R\ge3$ and $I>1$ imply $V_*>1$. Rearranging
\eqref{eq:auxnormfresh} yields
\[
 y_*^2-V_*^2=\frac{V_*^2-1}{I^2-1}>0.
\]
Thus $u=y_*-V_*$ is a positive integer. With $V+u=y_*$,
\eqref{eq:auxidentity} proves $A_a=0$. All the other residuals are zero
by construction, including the unchanged canonical index and transport
conditions. This completes both existence implications for $\Ptwo$.

\section{Exact degree and literal arithmetic cost}
\label{sec:degree}

All fifteen supplied unknowns have degree one; the four fixed ports have
degree zero. Degree refers to the polynomial after exact cancellation,
not an upper bound obtained by adding degrees at every raw circuit gate.
Two local identities expose the relevant cancellations. With
$Z=X+\gamma H$ and $D_0=ac+Z$,
\begin{align}
 A_0={}&q^6\bigl[r^2+\ell(2r-k)-1\bigr]
       +2q^3(\ell+r)Y+Y^2,\label{eq:firstcancel}\\
 D_0^2-\Delta c^2={}&2acZ+Z^2-Hc^2,\label{eq:maincancel}\\
 A_1={}&q^2\bigl[2acZ+Z^2-Hc^2-1\bigr]
             +2qD_0XH+(XH)^2.\label{eq:clearcancel}
\end{align}
These are exact integer polynomial identities. The cancellation-expanded
certificate circuits implement them literally.

The degrees of $a,H,c,\Delta,I,U,\widehat R,V$ are respectively
$2,2,2,4,5,6,3,5$. Equations \eqref{eq:firstcancel}--\eqref{eq:clearcancel}
then give the exact residual degrees
\[
 (\deg A_0,\deg A_1,\deg A_k,\deg A_t,\deg A_a,\deg A_s)
       =(11,9,4,3,12,10).
\]
Only $A_a^2$ contributes to total degree $24$. Its leading homogeneous
form is
\begin{equation}
 \operatorname{in}_{24}(\Ptwo)
 =4X^4Y^6f^8u^2\bigl[(\eta+\zeta)T+hX\bigr]^2.
 \label{eq:leader24}
\end{equation}
The coefficient of $X^6Y^6f^8u^2h^2$ is $4$, independently of the fixed
tuple. This proves \emph{exact} degree $24$, not merely degree at most $24$.

The supplied evaluator uses $90$ binary operations. An addition,
subtraction, or multiplication is one operation, including multiplication
by a small fixed numeral. Integer constants and references to earlier
gates are free leaves. Every power is implemented by multiplication.
Common subexpressions can be shared, and dead gates are pruned. This is
the same source-level metric used for the comparisons; it does not count
bit complexity or charge for the large coefficient values. The $79$-gate
shared residual evaluator gains six squarings and five additions for the
single sum of squares.

The cancellation-expanded circuit uses more gates and is only a degree
certificate. It is not substituted for the $90$-operation evaluator when
reporting cost. Both circuits compute the identical polynomial, as follows
from the local identities and the literal downstream wiring.

\section{Two fifteen-variable degree--cost tradeoffs}
\label{sec:tradeoffs}

\subsection{Degree 26 with 88 operations}

Keep the fifteen variables and all common expressions of
Section~\ref{sec:formula}, but use
\[
 V=f^2(cT-R)-c,\qquad N_k=k-hE-R,
 \qquad A_{kt}=N_kW-q.
\]
Define
\begin{equation}
 \Pthree=A_0^2+A_1^2+A_{kt}^2+A_a^2+A_s^2.
 \label{eq:p26}
\end{equation}
The first two residuals again recover $Y=sq^3$ and $X=wq$. Since
$W=qN_t$, the coupled residual is $q(N_kN_t-1)$ and gives
$N_kN_t=1$. The auxiliary reconstruction is unchanged. These source
factors satisfy the ordinary product equation \eqref{eq:oldproduct},
whose signed-unit proof supplies soundness. Canonical completeness gives
$N_k=N_t=1$ directly.

The residual degrees are $(11,9,7,13,10)$. With $q_0=3\beta j$, the
degree-$26$ leading form is
\[
 4X^4Y^4f^8u^2q_0^6(q_0-F)^2.
\]
It is nonzero for the actual positive $\beta$. The literal evaluator has
$88$ gates. Keeping the degree-four $R$ inside $V$ costs two degrees in
the final polynomial, while coupling the index and transport saves two
gates. This makes the role of the last index reduction transparent.

\subsection{Degree 30 with 79 operations}

Use positive unknowns
\[
 J,F,\alpha,z,f,h,i,T,w,Y,r,\eta,\zeta,u,\gamma.
\]
Here $J$ is unrestricted within $\Zp$, $q=\beta J+1$, and $X=wq$ is
again a polynomial abbreviation. Keep $Y$ directly supplied and use
$V=f^2(cT-R)-c$. Put
\[
 N_1=D_0^2-\Delta c^2,\quad N_k=k-hE-R,\quad
 N_t=(K_0+w)C+(q-F)-z(q-1).
\]
Then
\begin{equation}
 \Pone=A_0^2+(N_1N_kN_t-1)^2+A_a^2+A_s^2.
 \label{eq:p30}
\end{equation}
The first root alone forces $q^3\mid Y$. The paired residual makes
$N_1,N_k,N_t$ integer units. Since $\Delta=(a+2)^2-1$ is $0$ or $3$
modulo $4$, $N_1=-1$ is impossible modulo $4$. Thus $N_1=1$ and
$N_kN_t=1$, giving the same source product. No coprimality restriction
or main-root offset is needed. Completeness uses $r$ and $u$ as before
and keeps $\gamma=\gamma_{\ord}$.

The four residual degrees are $(12,15,15,10)$; the evaluator uses $79$
gates. For an explicit exact-degree certificate, put
\begin{align*}
 q_0&=\beta J,&X_0&=wq_0,&R_0&=q_0^3(q_0-F),\\
 L_1&=8\gamma X_0^2Y^3k,&L_k&=-hX_0Y-R_0,&
 L_t&=w(q_0-F-\alpha)-zq_0.
\end{align*}
The leading homogeneous form is
\[
 \operatorname{in}_{30}(\Pone)
       =(L_1L_kL_t)^2+(2X_0^2Y^2f^4uR_0)^2.
\]
It is nonzero. In particular the monomial
$w^4J^{10}Y^4f^8u^2F^2$ has coefficient $4\beta^{10}$, because the
first square contains no $u$. This point offers most of the degree gain
with four additional gates relative to the prior $44/75$ equation, or
three relative to the directly related $44/76$ root-gap equation.

\section{A twelve-variable companion of degree 629}
\label{sec:twelve}

The same integrality mechanism also improves the degree of a construction
with fewer supplied variables. The result in this section uses twelve strictly
positive integers and has exact degree $629$, with $114$ live binary arithmetic
operations. Its three-equation version has maximum residual degree $622$ and
$109$ operations, or $108$ after the established positive transport shift.
The comparison is with package~52's twelve-positive lean equation, of degree
$897$ and $108$ operations. Thus the new point reduces degree by approximately
$29.9\%$, at a cost of six operations. It does not reduce the number of supplied
variables, and no global optimality or literature-wide novelty is claimed.

All notation in this section is local. The four fixed ports are the same
complete authentic tuple $(\beta,K_0,M_C,M_F)$ used elsewhere in this report;
$M_F$ is already shifted. In particular $\beta=2^{590069}-1$. Canonical
completeness uses odd native exponent $N$, so its old coordinate satisfies
$J_{\rm old}\equiv1\pmod3$. We may therefore put $J_{\rm old}=3J+1$,
retaining completeness with positive $J$ by taking $N>1$. This changes the
parameterization, not the coefficient tuple. It gives $q\equiv2\pmod3$.
The point of this residue restriction is that the later divisor $q$ is
coprime to $3$.

\subsection{The fixed polynomial}

The twelve supplied coordinates are
\[
 J,F,\alpha,z,h,i,T,s,w,\eta,\zeta,t\in\mathbb Z_{>0}.
\]
Define the following polynomial abbreviations:
\begin{align*}
 j&=3J+1,&q&=\beta j+1,&C&=q-F-\alpha,\\
 X&=wq,&Y&=sq,&E&=XY,\\
 k&=\eta+\zeta,&c&=kY+\eta,&a&=E+Y,\\
 H&=4a+3,&\Delta&=a^2+H=(a+2)^2-1,\\
 \ell&=EkY,&A&=1+\ell(\ell+k),&G&=1+\Delta c^2,\\
 b_0&=X+ac,&&&
\end{align*}
\[
 R=\bigl((q-1)(q-F)+\alpha\bigr)(q^2-1)+(M_C+qM_F)j.
\]
\begin{align*}
 N_k&=k-hE-R,\\
 N_t&=(K_0+w)C+q-F-z(q-1).
\end{align*}
\begin{align*}
 I&=ic^2,&D&=1+\Delta I^2,&Q&=(\Delta I)^2,\\
 d_0&=Q-1,&v&=cT,&b&=c+RD,\\
 e&=Q(Dv^2+b^2)-1,&B_a&=2Qvb.
\end{align*}
The new scaled carrier is
\begin{align*}
 S&=q^2H,&L&=q^2(1+HIT),\\
 W&=q^2b_0+H\bigl[q^2(t+b+\ell)+s\bigr].
\end{align*}
Let
\begin{equation}
 \mathcal N=
 \prod_{\substack{g=\pm\sqrt G,\ \tau=\pm\sqrt A\\f=\pm\sqrt D}}
 \left[d_0(W-Lg-S\tau)^2-S^2e+S^2B_af\right].
 \label{eq:twelve-norm}
\end{equation}
Every formal sign is included. Consequently this is an integer polynomial;
no radicals or intermediate registers are additional unknowns. Formal signs
are not being identified with independent field automorphisms after
specialization. The equation is
\begin{equation}
 \boxed{P_{12}=N_kN_t(1+\mathcal N)-1=0.}
 \label{eq:twelve-final}
\end{equation}

For an explicit division-free evaluator, put
\begin{align*}
 d_A&=d_0S^2A,&U_Q&=S^2Q,\\
 P_H&=U_QDv^2,&U_H&=U_Qb^2,\\
 e_H&=P_H+U_H-S^2,&j_H&=4P_HU_H,\\
 \mathcal T&=2d_0(W^2+L^2G),&
 \kappa&=[d_0(W^2-L^2G)]^2,\\
 \mathcal A&=2(d_A+e_H),&
 \mathcal B&=(d_A-e_H)^2-j_H,\\
 a_r&=\mathcal T-\mathcal A,&b_r&=\mathcal B-\kappa,\\
 v_r&=16d_Aj_H-a_r(a_r\mathcal T+2b_r),&
 u_r&=b_r^2-a_r^2\kappa,\\
 \mathcal N&=\kappa v_r^2-u_r(\mathcal T v_r-u_r).
\end{align*}
This is the inherited quadratic-remainder identity for the formal norm.
The supplied JSON circuit gives every live gate and includes all
multiplications by fixed numerals. The system version compares the final
norm terms directly and imposes $N_k=N_t=1$.

\subsection{Bounds before recovering divisibility}

The predecessor norm identity gives
$\mathcal N\equiv(u_r-x_0v_r)^2\pmod4$, where
$x_0=d_0(W^2+L^2G)$. At a root of \eqref{eq:twelve-final}, all three factors
are integer units. The option $\mathcal N=-2$ is impossible modulo~$4$.
Thus $\mathcal N=0$ and $N_k=N_t=\varepsilon\in\{1,-1\}$.

The authentic native-mask inequalities applied to the transport unit give
\[
 C\ge0,\qquad q-F=C+\alpha>0,\qquad q\ge16,
 \qquad 3q<R<q^4-2.
\]
The boundary $C=0$ is not discarded at this stage. We cannot yet use
$q^3\mid Y$ or $E>R$: the present definition is only $Y=sq$. Instead,
\begin{equation}
 X,Y\ge q,\quad E\ge q^2,\quad
 k=hE+R+\varepsilon>R,\quad c>RY,\quad
 a+2<2E,\quad I\ge c^2.
 \label{eq:twelve-early}
\end{equation}
In bounds below, $g,f,\tau$ denote positive real square roots. Then
$g<2Ec$, $f<2EI$, and the inherited elementary expansion gives
\[
 \Delta c^2-A>
 k^2\bigl[EY(2Y^2-1)+Y^4\bigr]-1>0.
\]
In particular
\begin{equation}
 1<A<\Delta c^2\le\Delta I=\sqrt Q.
 \label{eq:twelve-exception-bound}
\end{equation}

The earlier lean construction used $E>R$ in its outer-root estimate. Here it
can be replaced by a consequence of $c>RY$. Since
$D<4E^2I^2$, $c<I^2$, and $d_0>\tfrac12E^4I^2$, we have
\[
 b+cf<(4R+1)E^2I^2,
 \qquad \frac{b+cf}{d_0}<\frac{9R}{E^2}.
\]
On the other hand, the exact identity
$(Ig)^2-(cf)^2=I^2-c^2$ implies
\[
 Ig-cf=\frac{I^2-c^2}{Ig+cf}>\frac{I}{8Ec}
       =\frac{ic}{8E}.
\]
Finally $icE>RYE=RXY^2>72R$ by \eqref{eq:twelve-early}. Therefore
\begin{equation}
 \frac{b+cf}{d_0}<\frac{9R}{E^2}<\frac{ic}{8E}<Ig-cf.
 \label{eq:twelve-new-bound}
\end{equation}
This is the extra estimate that makes divisibility recovery possible without
assuming the scale one is trying to recover.

\subsection{Integral roots force the missing scale}

For a coherent auxiliary root $y$ with $(f')^2=D$, the norm relation is
\[
 d_0y^2=Q(vf'-b)^2-1.
\]
The standard square estimate, also valid for complex $y$, gives
\[
 |y|<b+Tcf+T\frac{b+cf}{d_0}<b+ITg.
\]
Consequently any root $x=\tau'+y$ of the inner polynomial has
$|x|<\tau+b+ITg$. But the negative-main-root branch would give
\[
 \frac{W+Lg}{S}
 =b+ITg+\left(t+\frac{s}{q^2}\right)
   +\ell+\frac{b_0+g}{H}.
\]
This exceeds that upper bound: $t+s/q^2>0$,
$\tau<\ell+k$, and $(b_0+g)/H>c/3>k$, since $Y\ge16$.
Actual quadratic conjugation would send an irrational positive-main-root
branch to this forbidden branch. Hence $g$ is a positive integer and
$x=(W-Lg)/S$ is rational.

For completeness, the inherited inner-root reconstruction can be stated
explicitly. Define
\begin{align*}
 C_x&=d_0(x^2+A)-e,\\
 F(x)&=\bigl[C_x^2+DB_a^2-4d_0^2x^2A\bigr]^2-4C_x^2DB_a^2.
\end{align*}
The norm implies $F(x)=0$. If $C_x=0$, eliminating $x^2$ would give
\[
 d_0^2A^2-d_0eA+Q^2v^2b^2D=0,
\]
which forces $Q\mid A(A-1)$, contrary to
\eqref{eq:twelve-exception-bound}. Likewise $F(0)=0$ would force
$Q\mid(A-1)^2$, also impossible. Thus the denominators in
\begin{align*}
 U_x&=C_x^2+DB_a^2-4d_0^2x^2A,\\
 f'&=-\frac{U_x}{2C_xB_a},&
 \tau'&=\frac{C_x+B_af'}{2d_0x},&y&=x-\tau'
\end{align*}
are nonzero. The identity $F(x)=0$ gives
$(f')^2=D$, $(\tau')^2=A$, and the auxiliary Pell equation above.
The rational square roots $f',\tau'$ are integers. The inherited rational
Pell saturation lemma says that, for integers $m\ge2,V$ and rational $y$,
\[
 m^2V^2-(m^2-1)y^2=1\quad\Longrightarrow\quad y\in\mathbb Z.
\]
Applying it with $m=\Delta I$ and $V=vf'-b$ shows that $y$, and hence $x$,
is integral. The lemma's ordinary proof is preserved in the upstream proof
files; it is not a newly supplied proof-assistant certificate.

Reduce $W=Lg+Sx$ modulo $q^2$. This gives $q^2\mid Hs$.
But $Y=sq$ gives $H\equiv3\pmod q$, while $q\equiv2\pmod3$.
Thus $\gcd(q,H)=1$ and
\begin{equation}
 q^2\mid s,
 \qquad s_{\rm old}=s/q^2\in\mathbb Z_{>0},
 \qquad Y=s_{\rm old}q^3.
 \label{eq:twelve-divisibility}
\end{equation}

There is now a short soundness proof using the old equation directly.
Let $\mathcal N_{\rm lean}$ be package~52's unscaled twelve-variable lean norm,
on the same fixed tuple. In the localization at $q$ there is the exact identity
\begin{equation}
 \mathcal N=q^{32}
   \mathcal N_{\rm lean}\left(t+\frac{s}{q^2}\right),
 \qquad Y=(s/q^2)q^3.
 \label{eq:twelve-identity}
\end{equation}
Each of eight factors gains $q^4$. Its carrier degree is at most $16$, so
all denominators clear. By \eqref{eq:twelve-divisibility}, the mapped old
coordinates $s_{\rm old}$ and $t_{\rm old}=t+s_{\rm old}$ are positive
integers. The index and transport factors are unchanged. Thus a new root
produces a root of the established old positive equation. Its compiler and
K215 calibration recover a finite ZFC refutation.

\subsection{Completeness and coefficient-preserving equivalence}

Choose canonical normalized compiler witnesses with odd $N>1$. Keep the
retained coordinates, put $J=(J_{\rm old}-1)/3$, $s=s_{\rm old}q^2$, and set
\[
 t=\gamma+ITg+\tau+y-b-\ell-s_{\rm old}.
\]
All coordinates are integral. The auxiliary Pell equation gives
$y\ge b-vf$, and $Ig>cf$. Also
\[
 \tau-\ell=\frac{\ell k+1}{\tau+\ell}>\frac{k}{3},
 \qquad
 \frac{k}{3}>\frac{XY}{3}>\frac{Y}{q^3}=s_{\rm old}.
\]
Here $\ell>k$, $\tau<\ell+k$, $k>XY$, and $X\ge q\ge16$ are canonical
compiler inequalities. Therefore
\[
 t\ge\gamma+T(Ig-cf)+(\tau-\ell)-s_{\rm old}>0.
\]
A coherent factor in \eqref{eq:twelve-norm} vanishes, and $N_k=N_t=1$.
This proves both existence implications. Composing with the inherited
syntactic-refutation calibration yields, by ordinary mathematics formalizable
in ZFC,
\[
 \mathrm{ZFC}\vdash
 \bigl[\text{$P_{12}$ has no positive integer zero}\bigr]
 \longleftrightarrow\mathrm{Con}(\mathrm{ZFC}).
\]
No truth value of either equivalent assertion is proved. The separately
supplied system replaces the product by $N_k=N_t=1$ and $\mathcal N=0$.
Its shifted transport version uses $(K_0+w)C=F+z(q-1)$; the established
positive transport shift preserves existence after the scale has been
recovered.

\subsection{Exact degree and reproducible evidence}

Set the leading homogeneous forms
\begin{align*}
 q_0&=3\beta J,&X_0&=wq_0,&Y_0&=sq_0,&E_0&=X_0Y_0,\\
 k_0&=\eta+\zeta,&c_0&=k_0Y_0,&I_0&=ic_0^2,&\ell_0&=E_0c_0,\\
 H_0&=4E_0,&Q_0&=E_0^4I_0^2,&R_0&=q_0^3(q_0-F),\\
 B_0&=R_0E_0^2I_0^2,&V_0&=I_0T\ell_0.
\end{align*}
The respective degrees of $H_0,\ell_0,Q_0,B_0,V_0,c_0$ are
$4,7,30,26,15,3$. The exact leading cancellations in the unscaled lean
norm give the following eight factor degrees:
\[
 79,79,79,79,71,71,67,65.
\]
Their sum is $590$. Four opposite $f/g$ branches have leading form
$-\operatorname{sgn}(g)4H_0^2Q_0B_0V_0$; the two matched branches with
negative $\tau$ have leading form $4H_0^2Q_0B_0\ell_0$; the other two
matched branches have leading forms
\[
 H_0^2Q_0B_0c_0,
 \qquad
 H_0^2Q_0B_0\left(2t-k_0-\frac{iTc_0}{E_0}\right).
\]
Their product is
\begin{equation}
 2^{12}H_0^{16}Q_0^8B_0^8V_0^4\ell_0^2c_0
 \left(2t-k_0-\frac{iTc_0}{E_0}\right).
 \label{eq:twelve-leading}
\end{equation}
The displayed denominator cancels against preceding factors. The result is
a nonzero polynomial, since $t$ is independent. These are the same exact
lean-norm cancellations as in the bundled predecessor proof, now with
$\deg X=\deg Y=2$; in particular the cancellation uses
$(Ig)^2-(cf)^2=I^2-c^2$, not independent treatment of these square roots.

The shift $s/q^2$ in \eqref{eq:twelve-identity} has degree $-1$, below
$\deg t=1$, so it does not change \eqref{eq:twelve-leading}. The factor
$q^{32}$ adds $32$ degrees, giving $\deg\mathcal N=622$. The index factor
has nonzero degree-$5$ leader $-hE_0$, and the transport factor has nonzero
degree-$2$ leader $w(q_0-F-\alpha)-zq_0$. Hence
\[
 \deg P_{12}=622+5+2=629.
\]
This is uniform for every authentic positive $\beta$ satisfying the stated
coefficient and completeness interface.

The companion checker performs $60$ modular formal-norm multiplications,
$59$ valid modular checks of \eqref{eq:twelve-identity}, and six complete
univariate polynomial computations: three with formal-port fixtures and
three with the actual fixed tuple. All six give degrees $622$ and $629$.
It also checks all four compressed and uncompressed coefficient hashes and
all fourteen bundled source-file pins. These are finite algebra diagnostics;
the unbounded equivalence rests on the preceding argument and the inherited
compiler theorems. The files \texttt{twelve/upstream/LEAN\_PROOF.md},
\texttt{INNER\_RECOVERY\_PROOF.md}, \texttt{ELEVEN\_POSITIVE\_PROOF.md}, and
\texttt{ORDINARY\_BLANK\_PROOF.md} preserve those ordinary mathematical
boundaries explicitly.


\section{The fixed coefficients and their provenance}
\label{sec:coefficients}

All four new equations use the complete compact native tuple from source
44 \cite{source44}. The binding is
\[
 \beta=B-1,\quad K_0=D_C+BD_R,\quad M_C=M_C,\quad
 M_F=M_{F,\mathrm{native}}+B-1,
 \qquad B=2^{590069}.
\]
The last shift is part of the source interface. Replacing it by the native
high mask changes the polynomial and invalidates the stated theorem.
Nor can individual ports be selected from different historical packages.

For completeness, the following is a finite specification of the numerals.
Let $\rho=2^{67}$, $L=8807$, so $B=\rho^L$. The ordered list of $943$
literal Wang tiles is supplied in the archive. Its edge labels
$(w,n,e,s)$ have respective widths $8,9,8,9$. Set
\[
 v_t=w+2^8n+2^{17}e+2^{25}s+127\,2^{34},\qquad
 c_t=1+2^{10}(2^{41}-v_t).
\]
Let $\mathcal E$ consist of selector positions $0,\ldots,942$, west
copies $1023,\ldots,1031$, north copies $1013,\ldots,1022$, east copies
$1003,\ldots,1011$, south copies $983,\ldots,992$, dummy positions
$1063,1064$, and anchors $1095,2195,3296$.
For selectors use the coefficients $c_t$. The west payload bits are
$0,\ldots,7,34$; north bits $8,\ldots,16,35$; east bits
$17,\ldots,24,36$; south bits $25,\ldots,33,37$; the anchors are bits
$38,39,40$. A copy or anchor carrying bit $j$ receives coefficient
$2^{10+j}$. Both dummy coefficients are zero. Define
\begin{align*}
 \mu&=2^{51}-2,\\
 D_C&=\rho^{30}+\rho^{62}+\rho^{1100}+\rho^{2201}
       +\sum_{e\in\mathcal E}c_e
                  (\rho^{4393-e}+\rho^{5495-e}),\\
 D_R&=\rho^{42},\\
 Z_{\mathrm{pad}}&=(2^{65}-1)\sum_{e=8792}^{8804}\rho^e
                      +(2^{19}-1)\rho^{8805},\\
 M_C&=B-1-\sum_{\substack{e\in\mathcal E\\e\ne1}}\rho^e
                       -2\rho^{1063},\\
 M_{F,\mathrm{native}}&=4+\mu(\rho^{4393}+\rho^{5495})
       +\sum_{j=0}^{9}\rho^{1013+j}+\sum_{j=0}^{8}\rho^{1065+j}
       +\rho^{2195}+\rho^{3296}+Z_{\mathrm{pad}}.
\end{align*}
All summation limits and exponents are fixed numerals. The archive supplies
both the literal tile list and the literal $430$-action K215 transition
table from which the coefficient builder reconstructs it, without running
the machine. Thus the specification makes no reference to whether the
machine eventually halts.

\begin{table}[htbp]
\centering
\begin{tabular}{@{}lrrr@{}}
\toprule
Port & Bit length & Raw bytes & Gzip bytes\\
\midrule
$\beta$ & 590069 & 73759 & 108\\
$K_0$ & 592884 & 74111 & 8940\\
$M_C$ & 590069 & 73759 & 193\\
$M_F$ (shifted) & 590070 & 73759 & 204\\
\midrule
Total & & 295388 & 9445\\
\bottomrule
\end{tabular}
\caption{Materialized fixed ports. These are numeral-port sizes, not a
claim about the coefficients of a fully expanded polynomial.}
\end{table}

Each gzip payload contains an unsigned little-endian integer. The binding
records both compressed and raw SHA-256 digests. The principal circuit
aliases are $\texttt{K0}\mapsto\texttt{K}$ and
$\texttt{MF}\mapsto\texttt{MF\_source}$. The other two aliases are unchanged.
The twelve-variable circuits use the same map.

The copied adapter proves the complement-clause lemma, native alignment,
carry and population bounds, and the actual-rank padding controller. Its
parameters are $d=67\cdot8807=590069$ and $\lambda=581197$. In particular,
the controller supplies powers of an odd prime for both padded dimensions,
so the odd-$N$ argument used in the new proof is part of this exact tuple's
completeness theorem. The native low mask has residue $2$ modulo $4$ and
the canonical low-bit disjointness gives $C\equiv1\pmod4$, whence
$R\equiv3\pmod4$. These facts are not presumed for arbitrary integer ports.

The independent inherited coefficient auditor was rerun. It authenticated
the source files and checked $66{,}062$ native basis identities,
$38{,}663$ one-bit mutation diagnostics, and $87{,}380$ exhaustive small
complement models, as well as the concrete bounds and controller residues.
Finite audits verify that the implemented constants satisfy the finite
parts of the compiler interface; the unbounded semantic theorem remains
the preserved ordinary mathematical proof.

\section{Domains and what ``no solutions'' means}
\label{sec:domains}

The compact result is stated over strictly positive integers. This is a
standard Diophantine convention, but the domain must be written explicitly.
If $\mathbb N$ denotes $\{0,1,2,\ldots\}$, the polynomial
\[
 P_{\mathbb N}(x_1,\ldots,x_{15})
       =\Ptwo(1+x_1,\ldots,1+x_{15})
\]
has the same existence statement over $\mathbb N^{15}$. Its exact degree
is still $24$, since translation preserves the top homogeneous component.

For a completely unrestricted integer equation, let
\[
 Q_j=1+z_{j1}^2+z_{j2}^2+z_{j3}^2+z_{j4}^2,
 \qquad 1\le j\le15,
\]
and put $P_{\mathbb Z}=\Ptwo(Q_1,\ldots,Q_{15})$. Lagrange's four-square
theorem makes every positive integer attainable as a $Q_j$. Therefore
\[
 \neg\exists\mathbf z\in\mathbb Z^{60}\;P_{\mathbb Z}(\mathbf z)=0
 \quad\longleftrightarrow\quad\Con(\ZFC)
\]
is again provable in ZFC. The polynomial has exact degree $48$: the degree
is at most twice $24$, and setting three coordinates in every block to
zero leaves an injective substitution $x_j\mapsto z_{j1}^2$ in the nonzero
top form. This translation is supplied for domain clarity, not claimed as
a competitive unrestricted-integer frontier. In particular, this work
does not improve the supplied ten-integer-variable constructions.

\section{Reproducibility and limits of verification}
\label{sec:checks}

The archive is designed to be usable independently of the working checkout.
The main deliverables are this PDF and its LaTeX source, the literal JSON
circuits, Python builders and checkers, the complete coefficient payloads,
source provenance, and the inherited proof files. Python verification uses
only the standard library. The scripts locate files relative to themselves
and also work under \texttt{python3 -O}; correctness checks use explicit
exceptions rather than removable \texttt{assert} statements.

The main checks have different logical roles:
\begin{enumerate}
 \item \textbf{Exact algebra and degree.} Five local sparse polynomial
 identities prove the cancellation and index-reduction rewrites. Leading
 homogeneous forms are propagated exactly through the certificate circuits.
 A completely vanishing leading form causes failure, rather than silently
 accepting an uncertain degree. The resulting forms include
 \eqref{eq:leader24} and the two tradeoff leaders.
 \item \textbf{Literal circuit integrity.} All references are defined before
 use; every charged operator is binary $+,-,$ or $\times$; every advertised
 unknown and port is live; all operation counts are computed from retained
 gates. Builders regenerate the published circuits deterministically.
 \item \textbf{Independent transcription.} A separate checker imports no
 project builder or algebra implementation. It authenticates the actual
 coefficient payloads, compares independently transcribed formulas with the
 primary and certificate circuits in $7{,}344$ scalar residual/output checks,
 and performs $72$ full univariate finite-field polynomial comparisons.
 These rays also provide matching degree lower bounds at the actual ports.
 \item \textbf{Twelve-variable companion.} Its independent checker performs
 $60$ formal-norm multiplication checks, $59$ valid old-norm substitution
 checks, and six full polynomial-ray calculations, including three at the
 authentic tuple. It authenticates its fourteen inherited source pins.
 \item \textbf{Coefficient and source authentication.} The main algebra
 verifier authenticates $51$ copied source files. The independent coefficient
 auditor reconstructs the finite native constraints and controller checks.
 Hashes detect changed files; they do not prove the semantic content of a
 proof document.
\end{enumerate}

The independent mathematical reviews were produced by separate AI-assisted
review agents. They checked the new soundness and completeness reductions,
the positive margins, the coefficient interface, the twelve-variable
pre-divisibility bounds, and the degree calculations. These are internal
reviews, not external peer review or proof-assistant certification. Their
notes are retained so that a human reviewer can examine both the arguments
and their stated boundaries.

No test searches for a refutation of ZFC, assumes that a zero exists, or
infers consistency from a failure to find one. Evaluating randomly chosen
tuples does not establish an unbounded existence equivalence. That
equivalence rests on the displayed reductions and the inherited compiler
theorems. A full formalization would need to encode and check those
reductions against the exact bound constants and calibrated syntax.

\section{What is new, and what remains open}

The core new step is the reuse of monic norm roots to carry divisibility
conditions. Directly supplying $Y$ removes the repeated degree cost of
$sq^3$. A second carrier directly supplies $X$ after the first has made
its offset coprime to the modulus. A separate index equation then lowers
the auxiliary expression's degree. The squared-modulus quotient, the
ordinary strong-rank theorem, the root-gap technique itself, the compact
coefficient tuple, and the twelve-variable formal-norm evaluator are
credited inherited ingredients. The proof is an existence reduction and
may choose fresh auxiliary witnesses; it is not a same-coordinate
identity with the entire predecessor polynomial.

The twelve-variable result shows that the idea survives elimination of
explicit square-root witnesses, but it requires a new bound before the
missing scale is recovered. That order matters: using $E>R$ at that stage
would have assumed what the construction is supposed to prove. The
bound in Section~\ref{sec:twelve} instead uses $c>RY$ and $E\ge q^2$.

Several concrete directions remain. One can investigate whether another
existing monic relation can carry the remaining quotient conditions, apply
the lower-degree parent to selective witness elimination, and formalize the
two scale-recovery lemmas and the canonical positive-gap construction. The
literal circuits also leave room for an operation-count optimization that
does not change the mathematical polynomial. None of these further
improvements is included in the stated result.

There is no proof here that degree $24$ is minimal with fifteen positive
unknowns, or that these equations are simplest under an unspecified global
metric. What is established is a concrete new set of degree--cost points,
with fixed numerals, explicit equations, two-way reductions, and reproducible
arithmetic evidence, compared with a precisely identified source snapshot.

\appendix
\section{Reproduction commands and archive map}

From the unpacked archive root, the single command
\texttt{python3 verify\_all.py} runs the complete finite audit. The individual
commands are:
\begin{verbatim}
python3 code/build.py
python3 code/domain_translations.py
python3 code/verify.py
python3 code/independent_check.py
python3 twelve/build_p12.py
python3 twelve/check_p12.py
python3 coefficients/audit_coefficients.py
\end{verbatim}
The top-level runner executes these checks in order and records their
results. The LaTeX document can be rebuilt with:
\begin{verbatim}
cd report
pdflatex -interaction=nonstopmode -halt-on-error report.tex
pdflatex -interaction=nonstopmode -halt-on-error report.tex
\end{verbatim}
The second pass resolves cross-references and the contents. The twelve-variable
section is included from its own source file. The PDF was rendered and
visually inspected before packaging.

\begin{center}
\small
\begin{tabular}{@{}p{0.30\textwidth}p{0.63\textwidth}@{}}
\toprule
Directory or file & Purpose\\
\midrule
\texttt{report/} & Complete new paper, LaTeX source and rendered PDF.\\
\texttt{circuits/} & Three primary evaluators, their degree certificates,
 and explicit domain translations.\\
\texttt{code/} & Deterministic builders and exact/independent checkers.\\
\texttt{certificates/} & Machine-readable receipts and homogeneous forms.\\
\texttt{coefficients/} & Complete compact tuple, source data, native proof,
 builder, inherited independent auditor, and notices.\\
\texttt{twelve/} & Companion equation, systems, proof, checker, LaTeX section,
 and fourteen pinned upstream sources.\\
\texttt{upstream/} & Inherited ordinary kernel, Pell lemmas, comparison proofs,
 and semantic status documents.\\
\texttt{review/} & Explicitly labeled internal AI-assisted review notes.\\
\texttt{PROVENANCE.json} & Repository snapshot, exact source paths and digests.\\
\texttt{MANIFEST.sha256} & Hash inventory of the delivered files.\\
\bottomrule
\end{tabular}
\end{center}

\begin{thebibliography}{99}
\bibitem{source44}
ProveIt source 44, \emph{Complement-packed native clauses: a smaller complete
coefficient tuple}. Pinned source at the repository snapshot identified in
Section~\ref{sec:scope}; full proof and payloads reproduced in
\path{coefficients/}.\newline
\url{https://github.com/VladimirReshetnikov/ProveIt/tree/6d375f2d2deac0519e659bbce5b48087023c1430/Computability/DiophantineConsistency/con-zfc-diophantine-forms/44-b4-sparse/native}

\bibitem{source49}
ProveIt source 49, \emph{Support-aware Pell circuits}, section on the
fifteen-positive root-gap constructions. Preserved as
\path{upstream/OFFSET44_PRIOR.tex}.\newline
\url{https://github.com/VladimirReshetnikov/ProveIt/tree/6d375f2d2deac0519e659bbce5b48087023c1430/Computability/DiophantineConsistency/con-zfc-diophantine-forms/49-support-pell}

\bibitem{source47}
ProveIt source 47, \emph{Minimum-norm consistency polynomials}, supplement
with the degree-$44$, $75$-operation fifteen-positive equation.
Proof and literal circuit preserved under \path{upstream/MINIMA44_PRIOR*}.\newline
\url{https://github.com/VladimirReshetnikov/ProveIt/tree/6d375f2d2deac0519e659bbce5b48087023c1430/Computability/DiophantineConsistency/con-zfc-diophantine-forms/47-sparse-support/minima}

\bibitem{source51}
ProveIt source 51, \emph{A 66-operation consistency polynomial by a
squared-modulus quotient}. Preserved as
\path{upstream/SQUARED_MODULUS_PROOF.md}.\newline
\url{https://github.com/VladimirReshetnikov/ProveIt/tree/6d375f2d2deac0519e659bbce5b48087023c1430/Computability/DiophantineConsistency/con-zfc-diophantine-forms/51-sparse-arith}

\bibitem{source52}
ProveIt source 52, \emph{Sparse Pell norms and low-variable consistency
equations}. The twelve-positive lean proof, evaluator, and needed recovery
lemmas are preserved in \path{twelve/upstream/}.\newline
\url{https://github.com/VladimirReshetnikov/ProveIt/tree/6d375f2d2deac0519e659bbce5b48087023c1430/Computability/DiophantineConsistency/con-zfc-diophantine-forms/52-sparse-pell}

\bibitem{k215}
ProveIt, \emph{K215: 215 working states for ZFC consistency}, status at the
pinned snapshot. Preserved as \path{upstream/K215_README.md}.\newline
\url{https://github.com/VladimirReshetnikov/ProveIt/tree/6d375f2d2deac0519e659bbce5b48087023c1430/Computability/BusyBeaver/Research/ZFC/candidate-k215}

\bibitem{sun}
Zhi-Wei Sun, \emph{Ten unknowns for Hilbert's tenth problem over the
integers}, preprint, arXiv:2609.01594, version 2, 2026.\newline
\url{https://arxiv.org/html/2609.01594v2}

\bibitem{bayer2026}
Jonas Bayer, Marco David, Malte Ha{\ss}ler, Yuri Matiyasevich, and
Dierk Schleicher, \emph{Diophantine equations over $\mathbb Z$:
universal bounds and parallel formalization}, arXiv:2506.20909,
version 3, 2026.\newline
\url{https://arxiv.org/html/2506.20909v3}

\bibitem{bayer2025}
Jonas Bayer and Marco David, \emph{A Formal Proof of Complexity Bounds
on Diophantine Equations}, ITP 2025,
LIPIcs 352, article 3, 2025.\newline
\url{https://doi.org/10.4230/LIPIcs.ITP.2025.3}
\end{thebibliography}

\end{document}
